\chapter{The Action Layer}

The action layer is the boundary at which model decisions become effects on external
systems. A model selects an action; the action layer executes it, validates it, records
the result, and returns an observation. The design of this layer determines not only what
the agent can do, but how reliably and safely it does it.

A common but misleading simplification is to classify tools as deterministic. The
relevant distinction is between tools whose outputs are reproducible and whose side
effects are predictable, and those where neither holds. A schema validator is fully
reproducible. A database lookup is reproducible for a fixed state. A web search varies
because rankings change, pages are updated, and snippets are re-extracted. A sub-agent
wrapping another model introduces its own reasoning loop and may reach a different
conclusion on each invocation. The architecture must treat each tool according to its
actual reproducibility and side-effect profile.

\section{Schemas as Behavioural Contracts}

A tool schema is the primary mechanism by which the model learns what a tool does, when
it is appropriate to call it, and what constitutes a valid invocation. A well-designed
schema gives the tool an unambiguous name, a description that specifies both the
conditions under which the tool should be used and those under which it should not, typed
arguments with explicit constraints, and a structured output format.

Schema quality has a direct and measurable effect on tool selection fidelity. When two
tools have overlapping descriptions, the model resolves the ambiguity arbitrarily. When
argument names are vague, parameters are filled incorrectly. When output fields are
verbose or unstructured, the model may extract the wrong field or ignore the result
entirely. Research has confirmed that poor schema design is a primary driver of tool
hallucination---cases where the model selects the wrong tool, supplies malformed
parameters, or bypasses the tool entirely and fabricates an output that mimics what the
tool would have returned. The schema is therefore a behavioural control: its precision
determines the quality of the agent's decisions at the action boundary.

The following example illustrates a well-formed tool definition using the Strands
decorator pattern introduced in chapter~8. The name is unambiguous, the description
specifies both the purpose and the condition for use, arguments are typed with explicit
constraints, and the return type is structured:

\begin{lstlisting}[language=Python]
from strands import tool
from pydantic import BaseModel, Field

class JobStatusResult(BaseModel):
    job_id: str
    status: str  # "pending" | "running" | "completed" | "failed"
    progress_pct: float = Field(ge=0.0, le=100.0)
    error_message: str | None = None
    summary: str

@tool
def get_job_status(job_id: str) -> JobStatusResult:
    """
    Retrieve the current execution status of a processing job.
    Use this tool after triggering a job with run_job, and before
    retrieving results with get_results. Do not call this tool
    if the job_id is unknown or has not yet been submitted.
    """
    raw = _fetch_job_status(job_id)
    return JobStatusResult(
        job_id=job_id,
        status=raw["status"],
        progress_pct=raw.get("progress", 0.0),
        error_message=raw.get("error"),
        summary=(
            f"Job {job_id} is {raw['status']} "
            f"({raw.get('progress', 0):.0f}% complete)."
        ),
    )
\end{lstlisting}

Several design choices are worth making explicit. The description instructs the model
both when to call the tool and when to refrain---this negative condition is as important
as the positive one because it reduces spurious invocations. The return type is a typed
Pydantic model rather than a raw dictionary, which enforces the output contract at the
Python level and makes the schema visible to the framework. The \texttt{summary} field
gives the model a pre-formatted string it can incorporate directly into a response
without having to synthesise one from the structured fields.

\section{Read Tools and Write Tools}

Tools differ in their side effects, and that difference governs the level of autonomy
the agent is granted. Read tools---search, retrieval, inspection, calculation---carry no
persistent consequences and are appropriate for autonomous invocation. Write tools---those
that modify external state, send communications, or trigger processes---might require
explicit authorisation in the schema and, for consequential operations, a human approval
step before execution.

\section{Agent-as-Tool and the Opacity Problem}

When a subtask is sufficiently complex to warrant its own reasoning loop---a research
task that requires searching multiple sources, a code review that requires inspecting a
patch against a set of criteria, a planning task that decomposes a broad request into
ordered steps---it can be encapsulated as a sub-agent and exposed to the parent agent as
a tool. The parent agent supplies an input and receives an observation; the internal
execution is hidden.

This pattern introduces a structural opacity problem. The parent agent cannot inspect
the sub-agent's intermediate decisions unless the sub-agent exposes them through a shared
tracing infrastructure. Research on multi-agent failure modes has found that coordination
breakdowns---failures at precisely this boundary---represent the largest single category
of failure in production multi-agent systems. The consequence is that an agent-as-tool
pattern requires nested traces: the sub-agent's execution must appear as a child span
within the parent trace, with its tool calls, intermediate reasoning, and output all
visible at the same level of detail as any other tool invocation. The diagnostic value
of the trace collapses at the delegation boundary when this structure is absent.

\section{Summary}

The action layer exposes external capabilities through typed invocation boundaries. Each
tool is a contract: its schema governs how the model selects and invokes it, its read or
write classification governs how much autonomy is appropriate, and its output design
governs how well the model can reason from the result. When a tool is itself an agent,
the contract extends to require that the sub-agent's execution be fully traceable from
the parent context. The evaluation chapter builds directly on this structure: scores
attached to trajectory spans measure precisely whether each tool was selected correctly,
invoked with valid arguments, and its result incorporated appropriately into subsequent
reasoning.